\chapter{Conclusions}
\label{chap:conclusion}

\section{Summary}

This handbook documents the \textbf{Wu Fortran piecewise PV inversion
pipeline} \cite{Davis1992} applied to the 2025-01-08 California
blocking event, on a nine-level ERA5 grid with the perturbation split
into a \textbf{surface}, \textbf{lower} and \textbf{upper} piece:

\begin{enumerate}\setlength\itemsep{2pt}
  \item \textbf{Download / climatology / grid} (\texttt{shared\_steps/01--04}) ---
        ERA5 event hour, 30-year mean, and the Fortran \texttt{.grid} inputs.
  \item \textbf{Pass A/B} (\texttt{pvpialln}) --- boundary $\theta$,
        interior Ertel PV (850--200\,hPa), balanced $\psi$, for the mean
        and event states (step~05).
  \item \textbf{Pass C} (\texttt{qinvert21}/BALNC) --- total balanced
        reference (step~06).
  \item \textbf{Pass D} (\texttt{qinvertp21}/BALP) --- the three-piece
        perturbation inversion, \texttt{INLIN}=1 so the pieces sum to the
        total (step~07).
  \item \textbf{Diagnostics} --- Wu-vs-ERA5 PV, induced winds, PV
        advection, cross-sections, 3-D structure, closure (steps~08--11).
\end{enumerate}

\section{Key Findings}

\begin{enumerate}\setlength\itemsep{3pt}

  \item \textbf{Surface heating is a boundary condition, not interior
        PV.}  Interior Ertel PV exists only at $K{=}2..N_L{-}1$
        (850--200\,hPa); the surface (1000) and top (100) levels carry
        the boundary $\theta$.  By Bretherton's equivalence the warm
        surface $\theta'$ is a positive PV \emph{sheet} at the boundary,
        applied as an inhomogeneous Neumann condition.  We invert it as
        its own \textbf{surface piece} (\texttt{QLV}=$\{1\}$), isolating
        the surface thermal forcing.

  \item \textbf{The surface-piece $\psi'$ is predominantly cyclonic for
        this event} ($\sim$84\% of points have $\psi'<0$ at 1000\,hPa),
        because the domain-mean surface $\theta'$ is net warm
        ($\approx+1$\,K).  This is case-dependent, not a law: a net-cold
        surface anomaly would give an anticyclonic surface piece.

  \item \textbf{Wu $Q\to$ SI is the exact constant $10^{-8}$} when
        $\kappa=2/7$ (\Cref{chap:wu_coeffs}).  After the raw-$T$ input
        fix, $Q_{\text{Wu}}\times10^{-8}$ matches ERA5 native Ertel PV
        with median ratio $1.00$, $\corr=0.995$.  The former
        ``$\sim$3.5$\times$ residual'' and any empirical scaling are
        retired.

  \item \textbf{The 9-level (top-100) grid closes the piecewise sum
        better} ($R^2\!\approx\!0.44$ at 250\,hPa) than a uniform-top
        8-level (top-200) grid ($R^2\!\approx\!-0.2$).  The non-uniform
        top is accepted because all analysis of interest is at
        $\le 200$\,hPa.

  \item \textbf{Pass D converges (\texttt{INLIN}=1)}; the earlier
        ``\texttt{INLIN}=0 mandatory'' claim is obsolete.  Pass C's
        ``started diverging'' is a benign post-convergence heuristic.
\end{enumerate}

\begin{table}[H]
\centering\small
\caption{SI unit conventions for all pipeline outputs.}
\label{tab:si_units}
\begin{tabular}{lll}
\toprule
Quantity & Symbol & Units \\
\midrule
Streamfunction          & $\psi$          & m$^{2}$\,s$^{-1}$ \\
Geopotential            & $\Phi = gz$     & m$^{2}$\,s$^{-2}$ \\
Ertel PV                & $q$             & K\,m$^{2}$\,kg$^{-1}$\,s$^{-1}$ \\
PV advection tendency   & $-\mathbf{u}\cdot\nabla q$ & K\,m$^{2}$\,kg$^{-1}$\,s$^{-2}$ \\
Induced wind            & $u,v$           & m\,s$^{-1}$ \\
Pressure                & $p$             & hPa \\
\bottomrule
\end{tabular}
\end{table}

\section{Limitations and Future Work}

\begin{itemize}\setlength\itemsep{2pt}
  \item \textbf{Flat lower boundary.}  The solver treats 1000\,hPa as a
        rigid lid everywhere and ignores topography; over the western-US
        land surface the true lower boundary is well above 1000\,hPa.  A
        terrain-following / surface-pressure lower boundary
        ($p=\mathrm{hyam}\,p_0+\mathrm{hybm}\,p_s$) with the actual
        near-surface $\theta$ would be the principal physical upgrade.
        The boundary-$\theta$ Neumann condition itself is standard and
        correct; the GFDL/CSU PPVI tool uses the identical scheme (plus a
        regional ``surgery'' mask of the boundary $\theta$).
  \item \textbf{Coarse vertical/horizontal resolution} ($1.5^\circ$,
        9 levels) limits the static-stability estimate and the closure.
\end{itemize}
